\subsection*{4-domain model summary}
\label{sec:sim:summary}

Table~\ref{tab:results} summarizes the simulation results across all four conditions
($N = 500$ per condition, seed 2026). Figure~\ref{fig:gca} shows GCA factor
percentages and the precision--GCA co-variation across conditions.

\begin{table}[h!]
  \centering
  \caption{4-Domain simulation results across experimental conditions ($N=500$ per condition).
  Bold values highlight the key anti-correlation and GCA factor.}
  \label{tab:results}
  \begin{tabular}{lrrrrrr}
    \hline
    Condition & Precision & Meta & Tool & Learning & $r_{\text{meta,GCA}}$ & GCA\% \\
    \hline
    Normal         & 0.728 & 0.656 & 0.550 & 0.664 & \textbf{$-$0.658} & \textbf{83.5\%} \\
    Disrupted      & 0.302 & 0.639 & 0.226 & 0.565 & \textbf{$+$0.730} & 81.2\% \\
    Nurse (spatial)  & 0.737 & 0.655 & 0.425 & 0.783 & $-$0.682 & 79.7\% \\
    Forager (spatial) & 0.793 & 0.647 & 0.741 & 0.804 & $-$0.599 & 81.0\% \\
    \hline
  \end{tabular}
\end{table}

\subsection*{The GCA factor}
\label{sec:sim:gca}

When we extract the first principal component from the three precision-monotone domains (tool use,
caste-appropriate learning, GCA score), it explains 83.5\% of variance in the normal condition.
This is a stringent test: the GCA factor emerges purely from the shared precision substrate without
any additional factor structure imposed. The result validates the single-parameter architecture---if
multiple independent mechanisms underlay these domains, we would expect a much weaker first
component.

Across all four conditions, the GCA factor remains between 79.7\% and 83.5\%, demonstrating
that the single-precision structure is robust even when precision distributions differ markedly
between conditions.

\subsection*{The metacognition--intelligence anti-correlation}
\label{sec:sim:anticorr}

In the normal condition, individual-level correlations reveal a striking pattern:
\begin{itemize}
  \item Meta--GCA: $\Rcorr = -0.658$ (negative; more intelligent bees are worse at opt-out)
  \item Meta--Tool: $\Rcorr = -0.725$ (strongly negative)
  \item Tool--GCA: $\Rcorr = +0.696$ (positive; tool use and intelligence share a substrate)
\end{itemize}

This pattern is illustrated in Figure~\ref{fig:corrmat}. The anti-correlations emerge directly from
the model's architecture: agents with high precision perform better on monotone-precision tasks
(tool, GCA) but are past the optimum for metacognitive accuracy (Section~\ref{sec:analytical}).
The model thus predicts that metacognition and performance are \emph{functionally dissociated at the
individual level} despite arising from the same parameter.

This is empirically testable: correlating individual GCA scores (from the battery
in~\cite{AnimalCog2024}) with opt-out accuracy in the same honeybees should yield $r < -0.3$. No
such test has been conducted, making this the model's primary falsifiable prediction.

\subsection*{Caste-appropriate learning}
\label{sec:sim:caste}

Caste effects on learning are substantial. In the spatial task condition, foragers achieve a mean
accuracy of 0.804 and nurses achieve 0.783, compared to 0.664 for the mixed caste normal condition.
The advantage arises not from higher precision (all normal conditions have similar precision
distributions) but from the caste-match amplification in Domain 3. Disrupting this match
(e.g., assigning nurses to forager-style spatial tasks) reduces performance below even the
normal mixed condition.

\subsection*{Disruption effects and the sign flip}
\label{sec:sim:disruption}

Circadian disruption reduces mean precision from 0.728 to 0.302 ($-58\%$). This impairs all
precision-monotone domains substantially (tool use: $-59\%$, learning: $-15\%$). Metacognitive
accuracy shows a smaller, non-monotonic change ($-2.6\%$), consistent with the inverted-U
prediction: the disrupted precision distribution ($\bar{p} \approx 0.30$) falls below the optimum
($\pstar = 0.56$), partially restoring calibration.

The most striking finding is the \emph{sign flip} of the metacognition--GCA correlation:
from $r = -0.658$ in the normal condition to $r = +0.730$ under disruption (Figure~\ref{fig:disruption}).
Under disruption, precision is uniformly low and highly variable across agents. Any residual
precision benefits all domains simultaneously---including both GCA and metacognitive accuracy---so
the anti-correlation disappears and is replaced by a positive relationship. This sign flip is a
strong falsifiable prediction: it predicts that the correlation between metacognitive accuracy and
GCA should switch sign in circadian-disrupted bees.
